\chapter{Known Limitations and Future Engineering Roadmap}
\label{chap:limitations}

\section{The Imperative for Technical Honesty}
\label{sec:lim:honesty}

A foundational principle of the NumLang project is **Absolute Computational Honesty**. Rather than concealing performance regressions or framing engineering compromises as deliberate features, this chapter provides a candid, technically rigorous examination of the areas where competing compilers defeated NumLang in the Showdown benchmarks (\cref{chap:evaluation}), explains the precise micro-architectural reasons for these losses, and outlines the planned engineering remedies in Phases 67--80.

\section{Analysis of Benchmark Losses in the Showdown}
\label{sec:lim:losses}

\subsection{Loss 1: Linked-List Deforestation in \texttt{nrev} and \texttt{append3}}
In the double naive reverse (\texttt{nrev}) and triple list append (\texttt{append3}) benchmarks, Microsoft Visual C++ (\texttt{MSVC-O2}) defeated NumLang-SC:
\begin{itemize}
    \item \texttt{nrev}: MSVC-O2 ran in **$53.60\,\mu\text{s}$**; NumLang-SC took **$264.90\,\mu\text{s}$**.
    \item \texttt{append3}: MSVC-O2 ran in **$14.90\,\mu\text{s}$**; NumLang-SC took **$104.70\,\mu\text{s}$**.
\end{itemize}

\paragraph{Root Cause Diagnosis.}
Why did the supercompiler lose on canonical deforestation benchmarks?
\begin{enumerate}
    \item \textbf{Failure of Identity Inversion}: In \texttt{nrev}, mathematically $\mathtt{reverse}(\mathtt{reverse}(xs)) \equiv xs$. A truly optimal transformation would reduce the entire double reversal to an instantaneous identity pointer assignment ($O(1)$). However, NumLang's distillation engine operates by inductive unfold/fold matching: it folded the inner \texttt{reverse} into an accumulating loop and the outer \texttt{reverse} into a second loop, residualizing two sequential traversal passes over heap memory.
    \item \textbf{C Allocator Cache Locality}: The C/C++ baseline in MSVC allocated all linked-list nodes in a contiguous memory block, achieving near-perfect L1 hardware prefetching. NumLang's scoped arena allocator was fast, but traversing two full passes over memory degraded cache hit rates relative to MSVC's tightly optimized pointer loop.
\end{enumerate}

\paragraph{Engineering Remedy (Phase 67).}
Phase 67 will introduce **Structural Inversion Laws** into the distillation engine. By embedding an automated theorem-prover check that verifies when an involution satisfies $f(f(x)) = x$, the compiler will fold the dual traversal directly into a zero-cost identity assignment.

\subsection{Loss 2: Backend Vectorization Gap in \texttt{fib\_matrix}}
In the coupled Fibonacci matrix exponentiation benchmark (\texttt{fib\_matrix}):
\begin{itemize}
    \item Rustc-O ran in **$3.50\,\mu\text{s}$**; NumLang-SC took **$14.60\,\mu\text{s}$**.
\end{itemize}

\paragraph{Root Cause Diagnosis.}
Both compilers derived the exact same asymptotic reduction: collapsing the $N=1000$ sequential loop into an $O(\log N)$ binary matrix exponentiation performing 10 matrix multiplications.

The performance discrepancy stems entirely from low-level backend code generation:
\begin{itemize}
    \item Rustc targets LLVM, which applied auto-vectorization to the $2 \times 2$ matrix multiplication inner loop, executing it in packed 128-bit SSE2/AVX registers.
    \item NumLang-SC was evaluated using the Cranelift backend. While Cranelift generates code in sub-milliseconds, its loop vectorizer does not yet support automatic SLP vectorization of $2 \times 2$ scalar matrix products. Cranelift emitted 8 scalar load/multiply instructions, incurring higher cycle latency per matrix step.
\end{itemize}

\paragraph{Engineering Remedy (Phase 73).}
Phase 73 will configure the LLVM backend as the default code generator whenever the recurrence solver synthesizes companion matrix exponentiation kernels, directly matching LLVM's AVX2 vectorization quality.

\subsection{Loss 3: Fixed-Size Array Constant Folding in \texttt{map\_map}}
In \texttt{map\_map} (mapping a 20-element static array through increment and doubling):
\begin{itemize}
    \item Rustc-O evaluated in **$0\,\text{ns}$** (compile-time constant folded); NumLang-SC took **$10.70\,\mu\text{s}$**.
\end{itemize}

\paragraph{Root Cause Diagnosis.}
LLVM's scalar evolution pass recognized that the input array had a small, statically known length of 20 elements, unrolling the loop completely and evaluating all elements at compile time.

NumLang's supercompiler treats arrays as symbolic indexed buffers unless a discrete recurrence is detected. Because no arithmetic progression was present, the driver generated a clean, deforested scalar loop rather than fully unrolling the 20 iterations into static constants.

\paragraph{Engineering Remedy (Phase 68).}
Phase 68 will introduce **Bounded Array Peeling**: unrolling fixed-size array operations when the iteration count $N \le 32$ and array contents are statically known.

\subsection{Loss 4: Micro-Loop Strength in \texttt{hofstadter}}
In the Hofstadter mutual recurrence benchmark:
\begin{itemize}
    \item MSVC-O2 ran in **$300\,\text{ns}$**; NumLang-SC took **$12.90\,\mu\text{s}$**.
\end{itemize}

\paragraph{Root Cause Diagnosis.}
MSVC's global optimizer hoisted loop invariants and unrolled the mutual switch dispatch into raw register jumps. NumLang-SC closed the mutual recurrence, but retained a knot branch check at the loop boundary, incurring branch prediction overhead on every cycle.

\section{Scope and Boundaries of Formal Proofs}
\label{sec:lim:proof_scope}

While NumLang's mechanized Lean 4 formalization (\cref{chap:lean_semantics,chap:preservation}) represents a groundbreaking achievement, scientific rigor demands clear documentation of what has and has not been verified:

\begin{itemize}
    \item \textbf{Mechanized (Zero Axioms)}: The small-step operational semantics of SSA MIR, symbolic driving, anti-unification (MSG), Hamilton distillation folding, process-tree compaction ($\eta$-reduction), and constructive WQO termination.
    \item \textbf{Unverified Passes (Roadmap Phases 70--72)}:
    \begin{enumerate}
        \item Reynolds defunctionalization tag generation (\texttt{src/mir/defunctionalize.rs}).
        \item Bareiss Simplex integer linear programming for polyhedral loop schedules (\texttt{src/mir/supercompiler/polyhedral\_ilp.rs}).
        \item Tiered JIT on-stack replacement and runtime memory management (\texttt{src/runtime/}).
    \end{enumerate}
\end{itemize}

These unverified passes rely on runtime SMT translation validation (\cref{chap:validation}) and differential fuzzing (\cref{chap:testing}) rather than machine-checked Lean 4 kernel theorems.

\section{The Future Engineering Roadmap: Phases 67--80}
\label{sec:lim:roadmap}

\Cref{tab:roadmap_future} details the planned development roadmap for NumLang across Phases 67 through 80.

\begin{table}[htbp]
\centering
\small
\begin{tabular}{clp{9cm}}
\toprule
\textbf{Phase} & \textbf{Milestone Title} & \textbf{Technical Deliverables and Objectives} \\
\midrule
67 & Structural Involution Fusion & Discover and eliminate mathematical involutions ($f(f(x)) = x$) during distillation; eliminate double reverse. \\
68 & Statically-Bounded Array Peeling & Unroll small fixed-size arrays ($N \le 32$) into compile-time constants, matching Rustc on \texttt{map\_map}. \\
69 & GHC Build/Foldr Codata Parity & Integrate Church-encoded stream fusion combinators into codata driving to match GHC \texttt{vector} fusion. \\
70 & Mechanized Defunctionalization & Formalize Reynolds closure defunctionalization soundness in Lean 4. \\
71 & Mechanized Polyhedral Scheduling & Formalize the Bareiss Simplex integer linear programming schedule validity in Lean 4. \\
72 & Verified Tiered JIT Soundness & Formalize On-Stack Replacement (OSR) and frame deoptimization in Lean 4. \\
73 & Default LLVM Recurrence Backend & Automatically route companion matrix exponentiation kernels to LLVM with AVX2 vectorization. \\
74 & Native GPU Residualization & Synthesize CUDA and ROCm compute kernels directly from polyhedral loop nests. \\
75 & WebAssembly Backend & Emit WebAssembly System Interface (WASI) bytecode via Cranelift's native Wasm emitter. \\
76--80 & Distributed Supercompilation & Cloud-scale distributed process-tree exploration across multi-node clusters. \\
\bottomrule
\end{tabular}
\caption{The NumLang Long-Term Engineering Roadmap (Phases 67--80).}
\label{tab:roadmap_future}
\end{table}


\section{Microarchitectural Analysis of Cache Misses in Linked-List Benchmarks}
\label{sec:lim:microarch}

To diagnose why MSVC-O2 outperformed NumLang-SC on \texttt{nrev} ($53.60\,\mu\text{s}$ vs $264.90\,\mu\text{s}$) and \texttt{append3} ($14.90\,\mu\text{s}$ vs $104.70\,\mu\text{s}$), we captured hardware performance counter traces using Intel VTune / Windows Performance Toolkit:

\begin{table}[htbp]
\centering
\resizebox{\textwidth}{!}{%
\begin{tabular}{lcccc}
\toprule
\textbf{Compiler} & \textbf{L1 Data Cache Misses} & \textbf{L2 Cache Misses} & \textbf{LLC Misses} & \textbf{Branch Mispredictions} \\
\midrule
\textbf{MSVC-O2} & $1,420$ & $310$ & $42$ & $12$ \\
\textbf{Rustc-O} & $12,890$ & $4,120$ & $980$ & $84$ \\
\textbf{NumLang-Base} & $8,140$ & $2,800$ & $610$ & $54$ \\
\textbf{NumLang-SC} & $6,890$ & $2,100$ & $490$ & $46$ \\
\bottomrule
\end{tabular}%
}
\caption{Hardware Performance Counter Event Traces on \texttt{nrev} ($N=1000$).}
\label{tab:hw_counters_nrev}
\end{table}

The hardware counter analysis reveals:
\begin{enumerate}
    \item \textbf{Cache Locality Dominance}: MSVC's contiguous node pool achieved a $5\times$ reduction in L1 data cache misses compared to NumLang-SC.
    \item \textbf{Dual Memory Traversals}: Because NumLang-SC residualized the double reversal into two passes, it traversed the 1,000-node list twice in sequence, reloading memory cells across cache eviction boundaries.
\end{enumerate}

\section{Theoretical Blueprint for Phase 67: Involution Fusion}
\label{sec:lim:phase67_blueprint}

An involution is a function $f$ that is its own inverse:
\begin{equation}
f(f(x)) = x \quad \forall x
\end{equation}
In Phase 67, NumLang will formalize an **Involution Rewrite Rule** within the term interner:
\begin{definition}[Involution Term Equality]
If function $f$ is registered with the verified attribute \texttt{\#[involution]}, the term interner satisfies:
\begin{equation}
\mathtt{intern\_call}(f, [\mathtt{intern\_call}(f, [x])]) \implies x
\end{equation}
\end{definition}
With this theorem active, $\mathtt{nrev}(\mathtt{nrev}(xs))$ will collapse directly into $xs$ during the very first symbolic driving step, completely eliminating both loops and reducing runtime to an instantaneous $O(1)$ pointer copy ($< 10\,\text{ns}$).
